% !TEX root = ../../mixed-sklyanin-cubics.tex
% Current result: 2026-11-22 / v54.
\section{中心一次元による制限と三次までの余乗法}
\label{mix:sec:central-spherical-reconstruction}

$u$ の符号対合は $c\ne0$ の一つのファイバーを保たない。
ここでは対合を取らず、全ての $c$ に存在する中心超平面 $u=0$ への制限を使う。
この制限の twist は $\varrho_S^2$ であり、次数二の中心元 $v$ による
制限の twist $\varrho_S$ とは異なる。
その結果、余乗法の二次・三次の構造が残る。
以下は、この構造を比較すれば実際の残余圏を復元できることと、
四次以降には独立な比較条件が残らないことを証明する。

\subsection{特異係数でも使える proper な残余線束圏}

二ブロックの係数 $\alpha=(\alpha^{(1)},\alpha^{(2)})$ と $c\in\C$ に対し、
\[
 A_{\alpha,c}
 =\left(\mathcal B_{\alpha^{(1)}}\otimes_{\C[u]}
                   \mathcal B_{\alpha^{(2)}}\right)
       /(\Omega_1+100\Omega_2+cu^3),\qquad
 S=A_{\alpha,c}/(u)=B_0/(w_0)
\]
と置く。境界 $S$ と $\mathcal B=\mathcal K_S$ は $\alpha,c$ に依存しない。
$\mathcal C_{\alpha,c}$ を全ての $\OO_{A_{\alpha,c}}(j)$ が厚生成する dg 圏、
\[
 \mathcal A_{\alpha,c}
 =\langle\OO_{A_{\alpha,c}},\OO_{A_{\alpha,c}}(1),
                   \OO_{A_{\alpha,c}}(2),\OO_{A_{\alpha,c}}(3)\rangle^\perp
       \cap\mathcal C_{\alpha,c}
\]
をその残余部分とする。
これらは proper な冪等完備圏であり、
全 coherent 圏への一致は smooth 係数でのみ用いる。
可換側でも $G-u^2L+du^3$ に対し、同じ定義で残余線束圏を取る。
以下の $\varrho_A$ は
$L_{\OO_A}\circ(1)|_{\mathcal A_{\alpha,c}}$、
$\varrho_S$ は境界の同じ window 規約による rotation である。

\begin{lemma}[全係数環上の proper な線束模型]
\label{cmo:lem:central-universal-line-model}
$T=\C[\alpha^{(1)}_1,\ldots,\alpha^{(2)}_3,c]$ 上で上記の代数を作ると、
各内部次数は有限自由である。
線束圏と残余線束圏には $T$-proper な有限生成対象による dg 模型があり、
任意の係数変更と両立する。
境界は固定圏 $\mathcal B$ の $T$ への底変更である。
中心超平面制限、その左右随伴、およびその adjunction の全構造は
$T$ 上で定義される。
特に境界上の comonad は相対 perfect 核を持ち、
その低 arity の構造を境界の自由な mapping 群で表示した係数は $T$ の元である。
これは源の相対 smoothness を主張しない。
\end{lemma}

\begin{proof}
active 変数の次数を一、$u$ の次数を零とする濾過を使う。
二次関係の PBW 判定は係数に関する多項式恒等式であり、
命題\ref{cmo:prop:contraction}で用いた
$\beta\otimes1-1\otimes\beta=0$ は独立な全 $\alpha$ 上でも成立する。
従って周囲代数の随伴代数は $B_0\otimes_\C T[u]$ である。
中心三次の初項は $w_0$ であり、これは $B_0$ 上の非零因子なので、
商の随伴代数は $S\otimes_\C T[u]$ となる。
固定内部次数ではこの濾過は有限であり、その各商は有限自由 $T$-加群である。
従って商代数の各内部次数も有限自由で、同じ議論は任意の係数変更と両立する。
随伴代数の noetherian 性と有限な内部次数の濾過により
元の代数も左右 noetherian である。

命題\ref{cmo:prop:proper-contraction}の potential 双加群解消と
共有する中心差の解消は、同じ多項式係数で $T$ 上に存在する。
相対 Koszul 解消と Gorenstein 双対を使う同命題の証明は
PID 上の自由加群の判定を除いてそのまま適用でき、
その判定は上の有限自由性で置き換えられる。
従って線束間の Hom は
\[
 R\operatorname{Hom}(\OO(i),\OO(j))
 \simeq A_{j-i}\oplus(A_{i-j-4})^\vee[-5]
\]
という cohomology と任意底変更を持ち、$T$ 上 perfect である。
七本の線束 $\OO,\ldots,\OO(6)$ は ambient Koszul 解消から
全 twist を厚生成し、四本の例外線束に沿う有限 mutation が
残余の有限生成対象を与える。
その自己 Hom は $T$ 上 perfect である。

境界の線束は $u$ による長さ一の解消を持つため、
制限と inflation はこれらの相対線束模型間の関手である。
有限 mutation による $p_L,p_R$ と合成すれば左右の随伴を $T$ 上で構成できる。
単位・余単位およびその全 coherence はこの実 adjunction から取る。
境界は固定した smooth proper 圏の底変更であり、
得られた $FF^R$ は境界の有限生成対象を $T$-perfect な対象へ送る。
従って相対 smooth な境界の二側の対角を使えば、
その双加群は境界の積上も perfect である。
以下の分裂と正規化した $A_n$-tower は同じ相対核に適用できる。
そこで使う境界の mapping 複体は $\C$ 上のものの $T$ への底変更なので、
その自由 cohomology 群で測った構造係数は多項式であり、
開部分への局所化で生じる有理関数ではない。
\end{proof}

\begin{proposition}[中心超平面制限の球面性と生成]
\label{cmo:prop:central-linear-spherical}
全ての固定係数 $\alpha,c$ において、
中心超平面制限は dg 関手
\[
 F_{\alpha,c}
   =-\otimes_{A_{\alpha,c}}^{\mathbf L}S:
       \mathcal A_{\alpha,c}\longrightarrow\mathcal B
\]
を与え、左右の随伴を持つ。
その unshifted cotwist と twist は
\begin{equation}\label{cmo:eq:central-linear-periods}
 \mathsf C_F=\operatorname{Cone}(\operatorname{id}\to F^RF)
          \simeq\varrho_A^{-1}[1],\qquad
 \mathsf T_F=\operatorname{Cone}(FF^R\to\operatorname{id})
          \simeq\varrho_S^2.
\end{equation}
従って
\[
 \mathsf C_F^3\simeq[1],\qquad
 \mathsf T_F^3\simeq[4].
\]
標準の随伴比較射 $F^R\to\mathsf C_FF^L$ は同型であり、$F$ は spherical である。
さらに $F$ は conservative で、
\begin{equation}\label{cmo:eq:central-linear-generation}
 \mathcal B=\operatorname{thick}(F\mathcal A_{\alpha,c}),\qquad
 \mathcal A_{\alpha,c}=\operatorname{thick}(F^R\mathcal B)
\end{equation}
となる。特に、境界の split-generator $G_S$ に対して
$F^RG_S$ は $\mathcal A_{\alpha,c}$ の split-generator である。
以上は可換側の残余線束圏にも成立する。
\end{proposition}

\begin{proof}
$u$ は $A_{\alpha,c}$ 上正則である。
実際、周囲代数では $u$ が正則であり、
$w_{\alpha,c}\bmod u=w_0$ は $B_0$ 上正則である。
$ua=w_{\alpha,c}b$ を $u$ で還元すると $b=ub'$ となり、
$u$ を消去して $a=w_{\alpha,c}b'$ を得る。

包含を $a:\mathcal A_{\alpha,c}\hookrightarrow\mathcal C_{\alpha,c}$、
その随伴を $p_L\dashv a\dashv p_R$、
$S$ からの inflation を $i_*$ と書く。
境界は smooth で線束により生成される。
各 $i_*\OO_S(j)$ は
\[
 0\longrightarrow\OO_A(j-1)\xrightarrow{u}\OO_A(j)
       \longrightarrow i_*\OO_S(j)\longrightarrow0
\]
で解消されるので、$i_*$ は境界の全対象を $\mathcal C_{\alpha,c}$ へ送る。
制限・inflation の左右随伴も線束圏に制限でき、
\[
 F^R=p_Ri_*|_{\mathcal B},\qquad
 F^L=p_Li_*(1)[-1]|_{\mathcal B}
\]
となる。$M\in\mathcal A_{\alpha,c}$ の制限は
$[M(-1)\xrightarrow{u}M]$ で計算される。
$0\le r\le2$ について、この両項の $\OO_A(r)$ からの Hom は、
$\OO_A(r+1),\OO_A(r)$ から $M$ への Hom であり零である。
従って $FM\in\mathcal B$ である。

補題\ref{cmo:lem:mixed-line-serre}の線束圏上の Serre 関手は $(-4)[5]$ である。
単位三角へ $p_R$ を適用すると
\[
 \mathsf C_F M=p_R(M(-1))[1]\simeq\varrho_A^{-1}M[1].
\]
ここで inverse rotation は
$R_{\OO_A(-1)}\circ(-1)$ であり、
$\OO_A(-1)$ は Serre 双対によって $\mathcal A_{\alpha,c}^{\perp}$ に属する。
三次中心商の Orlov/window 比較における $(3)\simeq[2]$ は、
有限 mutation で定義した rotation 上の自然同型である。
両方向の rotation は全 twist の線束圏を保存するので、この周期も
$\mathcal A_{\alpha,c}$ に制限される。
これは源の smoothness を必要としない。

次に $Y=i_*N(1)[-1]$、$N\in\mathcal B$ とする。
$r=1,2,3$ では
$R\operatorname{Hom}(\OO_A(r),Y)
 =R\operatorname{Hom}(\OO_S(r-1),N)[-1]=0$ である。
従って $Y\to ap_LY$ の correction は $\OO_A$ の生成圏に入る。
$(-1)[1]$ を施すと $p_R$ がこの correction を殺す。
中心元の二項解消とその双対の評価により、全圏の随伴比較射は
$i_*N\to(({-1})[1])(i_*N(1)[-1])$ の恒等写像である。
従って残余側の標準射 $F^R\to\mathsf C_FF^L$ も同型になる。
これは命題\ref{cmo:prop:invariant-restriction-spherical}と同じ随伴射の計算で、
中心元の次数を一にしたものである。
\cite[Theorem~5.1]{CMOAnnoLogvinenkoSpherical}の判定が適用できる。

twist も随伴射を保って計算する。
四本の負の線束に関する right mutation を使い、
\[
 R\operatorname{Hom}(i_*N,\OO_A(r))
   =R\operatorname{Hom}(N,\OO_S(r+1))[-1]
\]
を適用すると、$r=-4,-3,-2$ では零である。
従って $ap_Ri_*N\to i_*N$ の correction は $\OO_A(-1)$ だけである。
全圏の余単位の cone は $N(-1)[2]$ なので、octahedron は
\[
 \mathsf T_FN\longrightarrow N(-1)[2]\longrightarrow E_N,\qquad
       E_N\in\langle\OO_S(-1)\rangle
\]
を与える。$\mathcal B$ は $\OO_S(-1)$ の左直交に属するため、
これは right mutation の三角である。従って
$\mathsf T_F=\varrho_S^{-1}[2]=\varrho_S^2$ となり、
式\eqref{cmo:eq:central-linear-periods}とその周期を得る。

最後に proper 性を使う。
$FM=0$ なら $\mathsf C_FM=M[1]$、従って $M[1]\simeq M[3]$ となり、
有界な自己 Hom により $M=0$ である。
余単位三角を反復した射 $N\to\mathsf T_F^{3m}N=N[4m]$ は
$m\gg0$ で零であり、その cone は $\operatorname{thick}(F\mathcal A_{\alpha,c})$
に属する。ゆえにその直和因子 $N$ も属する。
単位の fiber を $\mathsf H=\mathsf C_F[-1]=\varrho_A^{-1}$ とすると
$\mathsf H^3=[-2]$ である。同様に
$\mathsf H^{3m}M=M[-2m]\to M$ の cone を用いれば
第二の生成性を得る。
この証明は proper な線束圏だけで行っており、
特異係数での全 coherent 圏の proper 性は仮定していない。
\end{proof}

\subsection{分裂した underlying comonad と有限 arity 判定}

\begin{lemma}[一次 Bockstein の消滅]
\label{cmo:lem:central-linear-split-comonad}
命題\ref{cmo:prop:central-linear-spherical}の余単位三角で
$\nu:\operatorname{id}_{\mathcal B}\to\varrho_S^2$ は零であり、
その nullhomotopy の空間は可縮である。
従って、余単位を含む underlying kernel は
\begin{equation}\label{cmo:eq:central-linear-split-comonad}
 K=FF^R\simeq\mathbf1\oplus J,\qquad
 J=\varrho_S^2[-1],\qquad
 \epsilon:K\longrightarrow\mathbf1
\end{equation}
と一意的に分裂する。ここで一意性は
指定した $\varrho_S^2$ と三角に両立する分裂について述べている。
この分裂だけで余乗法が決まるとは主張しない。
\end{lemma}

\begin{proof}
全ての $\alpha,c$ について、次数付き代数として
$A_{\alpha,c}/(u^2)=S[u]/(u^2)$ である。
従って $S$-自由解消は同じ係数のまま $u^2=0$ 上の複体へ持ち上がる。
中心元 $u$ の Bockstein は零である。
補題\ref{cmo:lem:spherical-bockstein}の証明を次数一で行うと、
残余の $\nu$ はこの Bockstein の一意的な持上げである。
correction は $\OO_S(-1)$ の生成圏に入り、
$\mathcal B$ がその左直交であることによって、
各対象上の消滅だけでなく dg 自然変換の零性を得る。

分離比較により $\mathcal B\simeq\Ku(Z)$ とし、$Z=V(G)\subset\mathbb P^5$ とする。
diagonal sector 公式
\cite[Theorem~5.39]{NFBFKKernels}では、
$H^aR\operatorname{Hom}(\mathbf1,\varrho_S^b)$ の恒等 sector は
$a$ が偶数のとき $\operatorname{Jac}(G)_{b+3a/2}$ であり、
奇数のときは零である。
非恒等 sector の必要条件は $a$ が偶数かつ $2b+3a=6$ である。
従って $i\ge0$ に対して
$H^{-1-i}R\operatorname{Hom}(\mathbf1,\varrho_S^2)=0$ となる。
この消滅は nullhomotopy の空間の全 homotopy 群を消し、
指定した三角の分裂の一意性を与える。
\end{proof}

\begin{proposition}[二次と三次の構造で決まる全余乗法]
\label{cmo:prop:central-comonad-three-arity}
$\mathcal B=\Ku(Z)$、$Z=V(G)\subset\mathbb P^5$ を滑らかな三次四様体とし、
式\eqref{cmo:eq:central-linear-split-comonad}の counit arrow を固定する。
$\mathcal S_n$ をその上の counital $A_n$-comonad 構造の空間とする。
$n\ge2$ に対して
$\mathcal S_n\to\mathcal S_{n-1}$ の非空 fiber は離散空間であり、
既存の lift を原点に選んだその平行移動空間は
\begin{equation}\label{cmo:eq:central-comonad-finite-parameters}
 \begin{array}{c|c}
  n&\text{平行移動空間}\\ \hline
  2&\operatorname{Jac}(G)_1\simeq\C^6\\
  3&\operatorname{Jac}(G)_0\simeq\C\\
  n\ge4&0
 \end{array}
\end{equation}
である。特に、実際の二つの全 coherent comonad 構造が
$A_3$ 構造まで同じ連結成分に属するなら、
その比較は全 coherence を保つ同値へ一意的に延びる。
任意の二次・三次データの存在障害が消える、という主張はしていない。
\end{proposition}

\begin{proof}
補題\ref{cmo:lem:normalized-comonad-uniqueness}で確認した
固定 counit に関する $A_n$-tower
\cite[Theorem~3.3 and Lemma~3.11]{CMOWillumsgaardAssociativity}を使う。
正規化複体は
\[
 N_n=R\operatorname{Hom}(K,J^n)
     =R\operatorname{Hom}(\mathbf1,J^n)
           \oplus R\operatorname{Hom}(J,J^n)
\]
であり、非空 fiber の $i$ 次 homotopy 群は
$H^{2-n-i}N_n$、$i\ge0$ である。
ここで $J^n=\varrho_S^{2n}[-n]$ なので、
$H^qN_n$ の二つの恒等 sector の Jacobian 次数はそれぞれ
\[
 \frac{n+3q}{2},\qquad
 \frac{n-1+3q}{2}.
\]
第一項は $q-n$ が偶数のとき、第二項は $q-n+1$ が偶数のときだけ現れる。
$q=2-n-i$ を代入すると、第一項は $i$ が偶数のとき
$\operatorname{Jac}(G)_{3-n-3i/2}$、
第二項は $i$ が奇数のとき
$\operatorname{Jac}(G)_{(5-2n-3i)/2}$ である。
$n\ge2$ では、後者は全て零であり、
前者も $i=0,n=2,3$ 以外は零になる。
非恒等 sector の必要条件は
\[
 n+3q=6\quad\text{または}\quad n-1+3q=6
\]
であり、同じ代入によっていずれも不可能となる。
従って非空 fiber は離散で、式\eqref{cmo:eq:central-comonad-finite-parameters}を得る。

二次構造には分裂 kernel の primitive な余乗法を原点として選べる。
その正規化した差は
$\mathbf1\to J^2=\varrho_S^4[-2]\simeq\varrho_S$ で表される。
三次構造では、二次構造を同定した後の差が
$\mathbf1\to J^3[-1]=\varrho_S^6[-4]\simeq\mathbf1$ で表される。
これは零次 Jacobian 成分の scalar である。
ただし三次の原点は、同じ二次構造を持つ一つの実 lift を選んで定める。

$n\ge4$ の非空 fiber は可縮である。
実全構造から来る $A_3$ 構造の連結成分上では各段に lift が存在し、
path の持上げによりその成分の全ての点にも lift がある。
従って同じ $A_3$ 成分にある二つの実構造の比較は、各段で可縮な選択を持って延びる。
極限を取れば全 homotopy-coherent comonad の比較となる。
ここでは四次の存在障害の自動消滅を主張する必要はない。
\end{proof}

\begin{corollary}[三次までの比較からの実 Morita 核]
\label{cmo:cor:central-three-arity-kernel}
非可換側と可換側の中心超平面制限について、
境界の実 dg 同値 $E$ と rotation の同型を固定する。
これによって式\eqref{cmo:eq:central-linear-split-comonad}の underlying counit arrow を
同定した後、二次・三次の comonad 構造が同じ $A_3$ 成分に属すれば、
全ての固定係数で残余線束圏の dg Morita 同値と制限の intertwiner を構成できる。
両側の残余圏が smooth proper なら、その同値は実 perfect 双加群核で与えられる。
さらに補題\ref{cmo:lem:central-universal-line-model}の係数環上で
二次・三次のデータを大域的に同定すれば、
この構成は係数環に線型な相対 Morita 同値を与え、係数変更と両立する。
両側が相対 smooth proper である開部分上では、
相対比較核も二側の積上 perfect となる。
\end{corollary}

\begin{proof}
命題\ref{cmo:prop:central-comonad-three-arity}によって全 comonad 同型を構成する。
境界の生成対象 $G_S$ に対し、
命題\ref{cmo:prop:central-linear-spherical}から $F^RG_S$ は源の生成対象である。
従って実際の随伴による coKleisli 積を持つ代数
\[
 \Gamma_{\alpha,c}
 =R\operatorname{End}(F^RG_S)
 \simeq R\operatorname{Hom}(KG_S,G_S),\qquad
 g\star f=g\,K(f)\,\Delta_{G_S}
\]
は源の dg Morita 模型となる。
命題\ref{cmo:prop:spherical-comonad-model}の証明のうち、
生成対象と積を同定する部分は smoothness を用いていない。
全 comonad 比較はこの代数を積と高次構造を保って比較し、
graph 双加群と生成対象との tensor が実 Morita 同値を与える。
同じ比較は生成対象上の $F$ の作用にも両立するので、
有限 cone と直和因子へ延長して制限の intertwiner を得る。
smooth proper の場合には graph が体上 perfect であることと
両側の対角の smoothness によって双加群自体が perfect となる。
特異係数ではこの最後の perfect 性を主張せず、
残余線束圏の Morita 同値と、その帰結である smoothness の同値だけを用いる。

相対版については、境界の mapping 複体を係数環へ tensor すればよい。
式\eqref{cmo:eq:central-comonad-finite-parameters}の非零な二つの群は自由であり、
他の必要な群は零なので、二次・三次の大域的な同定は全 comonad の
大域的な同定へ同じ tower によって延びる。
また $T$-proper 性により
$R\operatorname{Hom}_T(M,M)$ は有界である。
従って $\mathsf H^{3m}M=M[-2m]\to M$ が大きい $m$ で零となる
生成性の証明は、$\C$ 上の有限次元性を
$T$ 上の perfect 性に置き換えてそのまま成立する。
境界の相対生成対象の右随伴像が相対源の生成対象となり、
coKleisli 代数の相対 graph 双加群から $T$-線型の実 Morita 同値を得る。
unit と counit はこの生成対象上の同型から全圏へ延びる。
以上の構成は相対 kernel の合成、有限 cone、直和因子および
一定の mapping 複体の係数変更で行われるため、任意の係数変更と両立する。
相対 smooth proper な開部分では graph の underlying $T$-複体は perfect であり、
二側の相対 smooth な対角を tensor して
相対 enveloping 代数上の perfect 性を得る。
\end{proof}
